\section{Consequences}\label{sec:consequences}

Theorem~\ref{thm:main}, together with the separate three-dimensional
theorem of Chen, Feng, Li, Xi, and Xu \cite[Theorem~1.1]{ChenEtAl2026},
gives Petty's projection-volume inequality and its ellipsoid equality
case in every dimension $n\ge3$. We now apply this combined conclusion
to lower-degree projection bodies, affine functional inequalities, and
isoperimetry in normed spaces. The implications and equality assertions
are specified separately in each setting.

\subsection{Lower-degree projection bodies}

For a convex body $L\subset\R^m$, use the standard intrinsic volumes
$V_j(L)$, normalized by the Steiner formula
\[
 \operatorname{vol}_m(L+tB_2^m)
 =\sum_{j=0}^m\kappa_{m-j}V_j(L)t^{m-j},
 \qquad t\ge0,\qquad \kappa_0=1.
\]
This is the normalization in
\cite[author version, Section~2, p.~6]{OrtegaMorenoSchuster2021};
on a hyperplane we use its induced Euclidean structure. For
$1\le i\le n-1$, the degree-$i$ projection body $\Pi_iK$ is defined by
\[
 h_{\Pi_iK}(u)=V_i(K|u^\perp),\qquad u\in S^{n-1},
\]
where $K|u^\perp=\operatorname{proj}_{u^\perp}K$. In particular,
$\Pi_{n-1}K=\Pi K$ because $V_{n-1}$ on $u^\perp$ is its
$(n-1)$-dimensional volume.

\begin{corollary}[Lutwak--Petty lower-degree inequalities]
\label{cor:lutwak-petty}
Let $n\ge3$ and $1\le i\le n-1$ be integers, and let
$K\subset\R^n$ be a convex body with nonempty interior. Then
\begin{equation}\label{eq:lutwak-petty}
\begin{aligned}
 \frac{V_{i+1}(\Pi_iK)}{V_{i+1}(K)^i}
 &\ge
 \frac{V_{i+1}(\Pi_iB_2^n)}{V_{i+1}(B_2^n)^i}\\
 &=
 \frac{\left(\displaystyle\binom{n-1}{i}
              \frac{\kappa_{n-1}}{\kappa_{n-i-1}}\right)^{i+1}}
      {\left(\displaystyle\binom{n}{i+1}
              \frac{\kappa_n}{\kappa_{n-i-1}}\right)^{i-1}}.
\end{aligned}
\end{equation}
The lower bound is sharp, since $K=B_2^n$ attains it.
\end{corollary}

\begin{proof}
Theorem~\ref{thm:main} and the separately proved three-dimensional
result \cite[Theorem~1.1]{ChenEtAl2026} establish Petty's conjecture for
every $n\ge3$. Lutwak's conditional result \cite{Lutwak1990}, in the
explicit formulation of Ortega-Moreno and Schuster
\cite[author version, Introduction, p.~2]{OrtegaMorenoSchuster2021},
states that Petty's conjecture implies the Euclidean-ball lower bound
for these ratios for every $1\le i\le n-1$. Applying that implication
gives the inequality. For its explicit value, the Steiner normalization
gives
\[
 V_j(B_2^m)=\binom{m}{j}\frac{\kappa_m}{\kappa_{m-j}},
 \qquad
 \Pi_iB_2^n=V_i(B_2^{n-1})B_2^n.
\]
The $(i+1)$-homogeneity of $V_{i+1}$ now gives the displayed ball
ratio. Only the inequality conclusion of the conditional result is used.
\end{proof}

At $i=n-1$, the displayed constant is
$\kappa_{n-1}^{\,n}\kappa_n^{\,2-n}$, exactly the constant in
Equation~\eqref{eq:main}. The case $i=1$ was already established by
Lutwak, as recalled in
\cite[author version, Introduction, p.~2]{OrtegaMorenoSchuster2021}.
Thus the additional range furnished here is $2\le i\le n-2$ for
$n\ge4$. No classification of equality cases is asserted for
Corollary~\ref{cor:lutwak-petty}.

\subsection{A sharp affine \texorpdfstring{$L^1$}{L1} Sobolev consequence}

For functions $f_1,\ldots,f_n\in C_c^1(\R^n)$, define the mixed
determinant-gradient integral
\[
 \mathcal D_n(f_1,\ldots,f_n)
 =\int_{(\R^n)^n}
 \left|\det\bigl(\nabla f_1(x_1),\ldots,\nabla f_n(x_n)\bigr)\right|
 \,dx_1\cdots dx_n.
\]
The determinant is the volume of the spanned parallelepiped, with no
factor $1/n!$. Under $f_i(x)\mapsto f_i(Ax+b)$ with $A$ invertible,
this integral and the product of the $L^{n/(n-1)}$ norms below both
scale by $|\det A|^{1-n}$.

\begin{corollary}[Mixed determinant-gradient and affine Sobolev inequalities]
\label{cor:affine-sobolev}
Let $n\ge3$ be an integer, put $q=n/(n-1)$, and let
$f_1,\ldots,f_n\in C_c^1(\R^n)$ be nonnegative. Then
\begin{equation}\label{eq:mixed-affine-sobolev}
 \mathcal D_n(f_1,\ldots,f_n)
 \ge n!\,\kappa_{n-1}^{\,n}\kappa_n^{\,2-n}
       \prod_{i=1}^n\norm{f_i}_q.
\end{equation}
In particular, every nonnegative $f\in C_c^1(\R^n)$ satisfies the
affine $L^1$ Sobolev inequality
\begin{equation}\label{eq:affine-sobolev}
 \mathcal D_n(f,\ldots,f)^{1/n}
 \ge (n!)^{1/n}\kappa_{n-1}\kappa_n^{(2-n)/n}\norm{f}_q.
\end{equation}
Both constants are sharp in this class, through smooth approximation
of ellipsoid indicators in $BV(\R^n)$; no equality assertion for nonzero
$C_c^1$ functions is made.
\end{corollary}

\begin{proof}
Write $c_n=\kappa_{n-1}^{\,n}\kappa_n^{\,2-n}$.
Haddad's geometric functional
$\widetilde I_1(K_1,\ldots,K_n)$ integrates the absolute determinant
against the surface-area measures of the convex bodies $K_i$.
He uses
\[
 h_{\Pi K}(u)=\frac12\int_{S^{n-1}}|\ip{u}{v}|\,dS_K(v),
\]
which is the projection-volume normalization above by Cauchy's
projection formula. His mixed-volume identity and the
Aleksandrov--Fenchel inequality give
\[
 \frac1{n!}\widetilde I_1(K_1,\ldots,K_n)
 =V(\Pi K_1,\ldots,\Pi K_n)
 \ge\prod_{i=1}^n|\Pi K_i|^{1/n}.
\]
Here $V$ denotes mixed volume, and Haddad's $\omega_m$ is our
unit-ball volume $\kappa_m$; see
\cite[arXiv version, Introduction, p.~5]{Haddad2020}.
Theorem~\ref{thm:main} for $n\ge4$ and the separately proved
three-dimensional case \cite[Theorem~1.1]{ChenEtAl2026} therefore give
\[
 \widetilde I_1(K_1,\ldots,K_n)
 \ge n!c_n\prod_{i=1}^n|K_i|^{(n-1)/n}.
\]

Haddad's Lemma~5.1 and the equivalence of Equations~(18)--(21)
\cite[arXiv version, Section~5, pp.~16--17]{Haddad2020}
transfer this $p=1$ geometric inequality to the functional inequality
for $\mathcal D_n$, with critical exponent $p^*=n/(n-1)=q$.
At this endpoint his normalized weak profile for $K$ is $\mathbf1_K$,
with the same surface-area measure as
$K$ and $\norm{\mathbf1_K}_q=|K|^{(n-1)/n}$; see
\cite[arXiv version, Equation~(9) and Definition~2.2, pp.~8--9]{Haddad2020}.
Consequently the functional and geometric sharp constants agree, as
also follows from the constant relations after his Equation~(21).
This proves Equation~\eqref{eq:mixed-affine-sobolev}; taking all
$f_i=f$ gives Equation~\eqref{eq:affine-sobolev}, the $p=1$ case of
his Equation~(22).

For sharpness, take nonnegative compactly supported smooth
mollifications $f_j$ of the indicator of an ellipsoid $E$. They converge to
$\mathbf1_E$ strictly in $BV(\R^n)$ and in $L^q$. Define finite measures
$\nu_j$ on the sphere by
\[
 \int\psi\,d\nu_j
 =\int_{\{\nabla f_j\ne0\}}
   \psi\!\left(-\frac{\nabla f_j}{|\nabla f_j|}\right)
   |\nabla f_j|\,dx,
 \qquad \psi\in C(S^{n-1}).
\]
Reshetnyak's continuity theorem
\cite[author version, Theorem~1.3, p.~3]{Spector2011} gives
$\nu_j\rightharpoonup S_E$; the minus sign selects the outward normal
in $D\mathbf1_E$. The product measures also converge weakly on the
compact product sphere. Since the absolute determinant is continuous
and homogeneous in each variable, this gives
$\mathcal D_n(f_j,\ldots,f_j)\to
\widetilde I_1(E,\ldots,E)=n!|\Pi E|$. Since
$\norm{\mathbf1_E}_q^n=|E|^{n-1}$, ellipsoid equality in Petty's
inequality gives the stated constants in the limit. This is the weak
$p=1$ interpretation of sharpness.
\end{proof}

To compare the diagonal statement with Zhang's affine Sobolev
inequality, let $f\ne0$ be as in Corollary~\ref{cor:affine-sobolev}
and define the origin-symmetric convex body $Z_f$ by
\[
 h_{Z_f}(u)=\frac12\int_{\R^n}|u\cdot\nabla f(x)|\,dx.
\]
This integral is positive for every unit vector $u$: otherwise $f$
would be constant on every line parallel to $u$, contradicting compact
support and $f\ne0$. Thus $Z_f$ has nonempty interior. Haddad's
zonoid identity gives $|Z_f|=\mathcal D_n(f,\ldots,f)/n!$, and his
Blaschke--Santal\'o comparison
\cite[arXiv version, Equations~(24)--(25), p.~17]{Haddad2020} yields
\[
 |Z_f^\circ|^{-1/n}
 \ge \left(\frac{\mathcal D_n(f,\ldots,f)}{n!\kappa_n^2}\right)^{1/n}
 \ge \frac{\kappa_{n-1}}{\kappa_n}\norm f_q.
\]
The resulting bound on $|Z_f^\circ|$ is Zhang's sharp affine Sobolev
inequality \cite{Zhang1999}. Thus Equation~\eqref{eq:affine-sobolev}
is a stronger affine $L^1$ inequality; the case $f=0$ is immediate.
This application uses only $p=1$ and gives no conclusion for
Haddad's $p>1$ problems.

\subsection{Holmes--Thompson isoperimetry}

\begin{corollary}[Holmes--Thompson isoperimetry]
\label{cor:holmes-thompson}
Let $n\ge3$, let $B=-B\subset\R^n$ be the unit ball of a real normed
space, and let $C\subset\R^n$ be a convex body. Write $B^\circ$ for the
polar unit ball and put $I_B=\Pi(B^\circ)/\kappa_{n-1}$. Normalize the
Holmes--Thompson volume and boundary area by
\[
 \mathcal V_B(C)=\frac{|B^\circ|}{\kappa_n}|C|,
 \qquad
 \mathcal A_B(\partial C)=nV(C[n-1],I_B),
\]
where $V(C[n-1],I_B)$ is mixed volume with $n-1$ copies of $C$,
normalized by $V(C[n])=|C|$. Then
\begin{equation}\label{eq:holmes-thompson}
 \mathcal A_B(\partial C)^n
 \ge n^n\kappa_n\,\mathcal V_B(C)^{n-1}.
\end{equation}
Equality holds if and only if $B$ is an ellipsoid and $C=x+tB$ for some
$x\in\R^n$ and $t>0$.
\end{corollary}

\begin{proof}
The fixed-norm isoperimetrix formula and the equivalence with symmetric
unpolarized Petty are classical; see Martini--Mustafaev
\cite[Definition~1(i), Equation~(2), Section~5, and Theorem~12]{MartiniMustafaev2008}.
Minkowski's mixed-volume inequality gives
\[
 \frac{\mathcal A_B(\partial C)^n}{\mathcal V_B(C)^{n-1}}
 \ge n^n\frac{\kappa_n^{n-1}}{\kappa_{n-1}^{n}}R_n(B^\circ)
 \ge n^n\kappa_n.
\]
The last step applies Theorem~\ref{thm:main} to $B^\circ$ for $n\ge4$,
and the separate three-dimensional result
\cite[Theorem~1.1]{ChenEtAl2026} for $n=3$. Equality in the first step
means that $C$ is homothetic to $I_B$, while equality in the last makes
$B^\circ$, hence $B$, an ellipsoid. Then $I_B$ is homothetic to $B$,
giving exactly the stated equality cases.
\end{proof}
